\documentclass[11pt]{article}
\usepackage[utf8]{inputenc}
\usepackage{hyperref}
\usepackage{geometry}
\usepackage{enumerate}
\usepackage{booktabs}

\geometry{a4paper, margin=1in}

\title{Recruiter Guide: My Projects Portfolio}
\author{Ali Ayhan G{\"u}nder}
\date{June 28, 2026}

\begin{document}

\maketitle

This document is a quick tour of the repository and where to find the most relevant work.

\noindent\textbf{Repository root:} \texttt{My-Projects/}

\tableofcontents

\newpage

\section{Fast path (25 minutes)}
If you only have a short amount of time, start with:
\begin{enumerate}
    \item \textbf{GENFIT (team full-stack SWE project)}: a fitness-themed platform built with documented setup.
    \\ \textbf{Folder:} \texttt{GENFIT\_SWE\_PROJECT\_CMPE451\_FINAL\_VERS\dot{I}ON/}
    \\ \textbf{Start here:} \texttt{.../bounswe2025group2-main/bounswe2025group2-main/README.md}
    
    \item \textbf{Operating Systems (C / performance)}: systems programming projects.
    \\ \textbf{Folder:} \texttt{OPERATING\_SYSTEM\_PROJECTS/}
    
    \item \textbf{Statistics / analytics (Python)}: data analysis scripts and reports.
    \\ \textbf{Folder:} \texttt{STATISTICS\_PROJECTS/}
\end{enumerate}

\section{How this repo is organized}
Projects are grouped mostly by course or topic. The top-level \texttt{README.md} acts as a clickable index.

\subsection{Main folders and what they contain}

\begin{itemize}
    \item \textbf{Software engineering (team)}
    \begin{itemize}
        \item \texttt{GENFIT\_SWE\_PROJECT\_CMPE451\_FINAL\_VERS\dot{I}ON/}: Final iteration of our team project.
        \item \texttt{GENFIT\_SWE\_PROJECT\_V1\_CMPE352/}: Earlier iteration of our team project.
    \end{itemize}
    
    \item \textbf{Systems \& low-level}
    \begin{itemize}
        \item \texttt{OPERATING\_SYSTEM\_PROJECTS/}: C projects; some are best built/run on Linux/WSL.
        \item \texttt{COMPUTER\_ORGANIZATION\_LABS/}: Course labs.
        \item \texttt{System Programming\_PROJECTS\_C,C++,ASSEMBLY/}: C/C++/assembly exercises.
    \end{itemize}

    \item \textbf{Distributed systems}
    \begin{itemize}
        \item \texttt{DISTRIBUTED\_SYSTEMS\_CMPE476/}: Sockets and RPC projects.
    \end{itemize}
    
    \item \textbf{Algorithms \& data structures}
    \begin{itemize}
        \item \texttt{ANALYSIS\_OF\_ALGORITHMS\_PROJECTS/}: Course projects on algorithm complexity and design.
        \item \texttt{DATA\_STRUCTURES\&ALGOR\dot{I}TMS/}: Fundamental structures and algorithms.
    \end{itemize}
    
    \item \textbf{Data}
    \begin{itemize}
        \item \texttt{DBMS\_SQL\_PROJECTS/}: SQL-focused database design projects.
        \item \texttt{STATISTICS\_PROJECTS/}: Python analytics, datasets, and report assets.
    \end{itemize}

    \item \textbf{System simulation \& modeling}
    \begin{itemize}
        \item \texttt{SYSTEM\_SIMULATION\_IE306\_projects/}: Discrete event simulation models (clinic, inventory, ferry simulations).
        \item \texttt{OPERATION\_RESEARCH\_PROJECTS\_IE310/}: Linear programming and optimization projects.
    \end{itemize}

    \item \textbf{Computational science \& simulations}
    \begin{itemize}
        \item \texttt{PARTICLE\_BASED\_SIMULATIONS\_CMPE\_49G/}: Simulation of physical particle movements and behavior.
    \end{itemize}

    \item \textbf{Signal processing}
    \begin{itemize}
        \item \texttt{SIGNAL\&SYSTEMS\_CMPE362\_Projects/}: MATLAB and Python signal processing labs.
    \end{itemize}

    \item \textbf{AI \& healthcare}
    \begin{itemize}
        \item \texttt{AI\_IN\_HEALTCARE\_CMPE49T/}: Deep learning applications on healthcare data.
    \end{itemize}
    
    \item \textbf{Languages / coursework}
    \begin{itemize}
        \item \texttt{JAVA\_OOP\_CMPE160/}: Java object-oriented programming projects (e.g. Ant Colony, Navigation).
        \item \texttt{PYTHON\_PROJECTS\_CMPE150/}: Introductory Python programming.
        \item \texttt{PRINCIPLES\_OF\_PROGRAMMING\_LANGUAGES\_CMPE 260-PROLOG/}: Logic programming and parsing in Prolog.
    \end{itemize}
    
    \item \textbf{HCI / UI}
    \begin{itemize}
        \item \texttt{HUMAN\_COMPUTER\_INTERACT\dot{I}ON\_PROJECTS\_CMPE496/}: UI/UX design reports and slides.
    \end{itemize}
    
    \item \textbf{Entrepreneurship / business}
    \begin{itemize}
        \item \texttt{Entrepreneurship\_PROJECT\&STARTUP\_AD432/}: Business plans, decks, and market research.
    \end{itemize}
\end{itemize}

\section{How to review quickly (what to look for)}

\subsection{Code-first projects}
\begin{itemize}
    \item Look for a project-level \texttt{README} with setup steps.
    \item For C/C++ projects, check for a \texttt{Makefile} and a \texttt{src/} folder.
    \item For Python projects, check for \texttt{requirements.txt} and runnable scripts (e.g., \texttt{task1.py}).
\end{itemize}

\subsection{Report-first projects}
\begin{itemize}
    \item Some folders include PDFs (reports/specs). Opening the PDF first can be the fastest way to understand the goal and results.
\end{itemize}

\section{Optional: running the highlight project (GENFIT)}
The GENFIT repository contains detailed Docker instructions in its own README. In general terms, the flow is:
\begin{enumerate}
    \item Install Docker + Docker Compose.
    \item Follow the \texttt{docker-compose} commands in the project README.
\end{enumerate}

\section{Contact}
You can reach me via email or GitHub for questions regarding these projects.
\\ \textbf{GitHub:} \href{https://github.com/Ayhan-coder}{github.com/Ayhan-coder}

\end{document}
